\documentclass[11pt,a4paper]{ctexart}
\usepackage[margin=2.4cm]{geometry}
\usepackage{amsmath,amssymb}
\usepackage{fancyhdr}
\usepackage{enumitem}
\linespread{1.18}
\setlength{\headheight}{15pt}
\setlength{\parskip}{0pt}

\pagestyle{fancy}
\fancyhf{}
\fancyhead[L]{\small Team 2200289}
\fancyhead[R]{\small \thepage}
\renewcommand{\headrulewidth}{0.4pt}
\fancypagestyle{plain}{\fancyhf{}\renewcommand{\headrulewidth}{0pt}}

\ctexset{
  section/format = \Large\bfseries\raggedright,
  subsection/format = \large\bfseries\raggedright,
  subsubsection/format = \normalsize\bfseries\raggedright
}
\setcounter{secnumdepth}{3}

\begin{document}

%====================== 摘要页 ======================
\thispagestyle{empty}
\begin{center}
{\small 2022 MCM/ICM\quad 选题 A\quad 队伍控制号 2200289\quad 摘要页（Summary Sheet）}\\[18pt]
{\LARGE\bfseries 自行车赛中的博弈论}\\[8pt]
{\large （Game Theory in Cycling）}\\[10pt]
\end{center}

\begin{center}
{\large\bfseries 摘\quad 要}
\end{center}

\vspace{6pt}
车手的策略对比赛结果有着巨大的影响。本文分析了车手功率输出的生理学模型与动力学模型，并在此基础上借助模拟退火算法（Simulated Annealing Algorithm）求得了整个赛道上的最优功率输出策略。随后，我们对可能的影响因素进行了敏感性分析，并把模型推广到了团体计时赛。

在模型一中，我们基于生理学建立了描述车手功率上限与体力的模型。通过对文献中的数据进行拟合，我们得到了刻画车手身体能力的定量参数。接着，我们建立了描述功率输出与持续时间之间关系的模型，并利用该模型给出了三类不同车手的功率曲线。

在模型二中，赛道被划分为 1164 个路段。我们首先对功率输出、动力学模型与天气条件进行分析，从而定量求解赛道上每一个路段的能量变化、速度变化以及所需时间。我们以缩短总比赛时间为优化目标，采用模拟退火算法对整条赛道上的功率分配进行优化。随后，我们把该模型应用于两条真实赛道和一条自行设计的赛道，并相应给出了最优功率分配方案。所有结果都得到了恰当的说明。

在敏感性分析中，我们首先通过在两类不同赛道中改变风速与风向这两个参数，考察了它们对模型二计算结果的影响。我们发现，赛道越接近环形，模型对天气条件的敏感程度就越低。

在分析模型优化结果对车手功率分配偏差的敏感性时，我们发现总时间与扰动方差呈正相关，并且对其相当敏感。随后，我们依次对每一个路段进行敏感性分析，以确定哪些路段属于关键路段，并反向求解出关键路段上可能的分段计时区间。

在模型推广部分，我们首先研究了把车队视为一个整体的简化情形。我们建立了描述全体排列方式下每一位车手所受空气阻力的模型，随后研究了在特定位置牺牲两名破风手的复杂模型。在模型二的基础上，我们改变目标函数，并把抛弃点作为决策变量引入，从而把模型推广到团体计时赛。

最后，我们针对某条具体赛道给出了一份比赛指导，用以介绍我们的模型并演示其使用方法。

\vspace{6pt}
\noindent\textbf{关键词：}生理学；动力学；模拟退火算法

\newpage
\pagestyle{fancy}

%%====================== 正文 ======================
\section{引言}

\subsection{问题背景}

自行车运动是当今最受欢迎的现代竞技体育项目之一。按照比赛形式划分，自行车公路赛共有三种类型：绕圈赛（criterium）、团体计时赛（team time trial）与个人计时赛（individual time trial）。这三种形式虽然同属公路骑行，但对车手能力的要求各有侧重：绕圈赛赛程较短、弯道密集、冲刺频繁，对爆发力与技战术位置感要求极高；团体计时赛以车队为单位统一计时，强调队形排列、轮流破风与整体节奏的一致性；个人计时赛则以车手独自对抗空气阻力为核心，考验车手在长距离上对自身功率输出的精确控制能力。在自行车比赛中，影响比赛结果的因素众多，包括运动员的能力、天气条件、赛道本身以及比赛策略。正因为这些因素相互耦合，与那些主要依赖运动员高爆发力的运动项目相比，在自行车运动中，依据特定运动员与特定赛道制定科学策略的重要性显得更为突出。

不同类型的运动员具有不同的身体特征，这一点不仅体现在产生大功率的能力上，也体现在功率能够维持多久上。爆发力强但耐力不足的运动员往往难以取得最佳成绩，反之亦然。一名车手在赛道上每一次功率输出的抬高，都会以消耗体内有限的能量储备为代价，因此“快”并不等于“全程都用力”，而是等于“在正确的位置用正确的力”。对运动员在运动过程中的身体变化进行数学建模，可以帮助教练制定最优策略，从而在给定运动员身体能力的前提下，使完成赛道所需的时间最短。科学的竞赛策略不仅能够帮助顶尖运动员打破纪录，对于自行车爱好者制定个人计划、节省体力而言同样具有意义。

\subsection{问题重述}

结合问题陈述中给出的背景信息与限制条件，我们需要建立一个对不同运动员都具有普适性的模型，并利用该模型完成以下任务：

\begin{itemize}
\item 给出两位不同性别的典型车手的功率曲线的定义。这一要求意味着模型不能只适用于某一位具体的运动员，而应当由若干可测量的生理参数决定，从而在不同性别、不同类型的车手之间实现参数化的切换。
\item 把你的模型应用于各种计时赛赛道。这要求模型能够读取赛道的地形信息，并在坡度、弯道等条件变化时依然给出合理的功率分配方案。
\item 研究天气条件对模型的影响，并对其作敏感性分析。这要求我们在模型中显式地引入风速与风向，并考察它们的变化如何改变最优解。
\item 研究车手对策略的偏离所带来的影响，并对其作敏感性分析。这要求我们量化“执行不到位”对总成绩的惩罚，并据此识别出对执行精度要求最高的路段。
\item 把模型推广到每队六名队员的团体计时赛的最优策略。这要求我们在单人模型的基础上引入队形排列、破风效果以及队员分工。
\item 为车队经理（Directeur Sportif）设计一份两页的自行车比赛指导，其中包含方向要点概述与模型概要。这要求我们把数学结论转写为可直接执行的赛场指令。
\end{itemize}

\subsection{我们的工作}

本题要求我们对车手的功率进行数学建模，并利用我们的模型设计最优的比赛策略。因此，我们的工作包括以下内容。本文的整体结构可以概括为一条自底向上的链条：第一层是基于生理学的功率上限模型，它回答“这名车手在给定时间内最多能以多大功率骑行”；第二层是包含动力学与天气条件的赛道推进模型，它回答“给定每一个路段的功率，车手会以多快的速度通过、需要多长时间”；第三层是优化层，它用模拟退火算法在全部路段上搜索使总时间最小的功率分配；第四层是敏感性分析层，它分别在天气条件与功率执行偏差两个方向上检验结论的稳健性；第五层是推广层，它把模型从个人计时赛扩展到团体计时赛；最后一层是应用层，它把上述结果整理成面向车队经理的比赛指导。

\section{假设与论证}

为简化问题，我们作出如下假设，并对每一条假设给出恰当的论证。这些假设的作用在于把复杂的实际问题化为可用微分方程与优化算法处理的形式，同时保证模型的主要机理不被破坏。

\begin{itemize}
\item \textbf{假设 1：}车手的体力始终在恢复，且恢复速率为常数。

\textbf{论证：}恢复速率是衡量有氧能力的指标，它与运动员的恢复能力相关。对同一名运动员而言，在整个比赛过程中恢复速率可以视为常数。也就是说，我们把车手在低于临界功率状态下骑行时无氧能量的回补过程近似为线性过程，其斜率不随赛程的推进而发生系统性变化。

\item \textbf{假设 2：}车手能够输出的最大瞬时功率与身体剩余的能量有关。

\textbf{论证：}当能量尚未被消耗时，人体能够爆发出最大功率；而当能量耗尽时，人体无法产生很大的功率。因此，假定车手的剩余能量决定了其运动表现的上限是合理的。这一假设使得功率上限可以写成剩余无氧能量的线性函数，从而把生理学描述转写为便于求解的代数关系。

\item \textbf{假设 3：}风向与车手的运动方向平行。

\textbf{论证：}根据流体动力学，当空气以一定速度撞击障碍物时，气流会沿其表面流动，从而与车手的运动方向平行，更何况在比赛赛道中偏斜角相当小（小于 22 度）。此外，由于地形复杂，对气流进行精确模拟难以实现，而这也不是我们研究的重点。因此，把风的作用简化为沿运动方向的一维作用，既能保留风阻的主要效应，又能使模型保持可解。

\item \textbf{假设 4：}车队中的每一名成员都具有相同的身体能力。

\textbf{论证：}在实际中，运动员之间身体能力的细微差异不可避免，而把它们纳入数学模型并不可行。因此，为简化问题并便于建模，我们假定团体比赛中的每一名运动员都具有相同的功率曲线。这样，团体计时赛的优化便可集中体现在队形排列与人员分工上，而不至于被个体差异所淹没。

\item \textbf{假设 5：}自行车队的队形变换在瞬间完成。

\textbf{论证：}车手完成队形变换只需要几秒钟，其间消耗的能量相对于整场比赛的能耗可以忽略不计。因此，我们可以把队形变换视为在某一位置处发生的瞬时事件，而不必为其单独建立过渡过程。

\item \textbf{假设 6：}在团体计时赛中，车手之间保持恒定的安全距离。

\textbf{论证：}为了在保证安全的同时把风阻降到最低，车手之间应当保持一个安全距离。考虑到职业自行车手的技术水平以及极少数剧烈的加速与减速路段，我们假定车手几乎在任何时刻都能够保持该距离。由此可以进一步认为，同一车队内各车手的速度与加速度始终一致。

\item \textbf{假设 7：}本研究中所使用的数据是准确的。

\textbf{论证：}我们假定所收集的车手数据是准确的，从而能够在此基础上建立合理的数学模型。若数据本身存在系统误差，则模型参数的拟合结果以及后续的优化结论都会随之产生偏移，因此这一假设是全部定量结论得以成立的前提。
\end{itemize}

\section{数据}

我们收集的数据包括赛道信息、若干组自行车运动员的成绩数据与身体数据，以及车手的动力学参数。这些数据分别服务于模型的不同层次：赛道信息用于确定每一个路段的坡度与曲率，运动员的成绩与身体数据用于标定功率曲线中的生理参数，而动力学参数（如阻力系数、迎风面积、摩擦系数、传动效率）则用于把功率换算为速度与位置的变化。

我们收集了 2021 年奥运会个人计时赛赛道与 2021 年 UCI 世界锦标赛个人计时赛赛道的等高线地形图。我们使用 Python 处理赛道的图像文件，以得到赛道上每一个路段的坡度。以 2021 年奥运会个人计时赛赛道为例，我们通过 Python 读取像素点，把赛道划分为一千多个节点，从而获得每一个节点的水平距离与垂直距离，以供进一步研究。水平距离与垂直距离之比即给出该路段的倾角，而节点的密集程度决定了后续数值积分的精度：节点越密，坡度变化被刻画得越细致，速度与时间的数值解也越接近真实情形。

\section{符号说明}

本文所用符号及其含义与单位如下所列。为便于阅读，凡在正文中首次出现的符号，均在此处给出统一的界定，后续各节不再重复说明。

\begin{itemize}
\item $P$：瞬时功率，即车手在某一时刻实际输出的功率，单位 W。
\item $P_c$：有氧运动的临界功率，即车手可以无限期维持而不消耗无氧能量的功率水平，单位 W。
\item $w$：无氧能量，即车手体内可供无氧代谢动用的能量余额，单位 J。
\item $AWC$：无氧工作能力（Anaerobic Work Capacity），即高于临界功率工作时可动用的无氧能量总量，单位 J。
\item $T$：以功率 $P$ 持续的时间，单位 s。
\item $P_{lim}$：车手的功率上限，即受剩余无氧能量约束时所能输出的最大功率，单位 W。
\item $\alpha$：$P_{lim}$ 中 $w$ 的系数，表征无氧能量对最大功率的贡献强度，单位 $\mathrm{s}^{-1}$。
\item $\rho$：空气密度，单位 $\mathrm{kg/m^3}$。
\item $S_{eff}$：有效迎风面积，即阻力系数与迎风面积的乘积，单位 $\mathrm{m^2}$。
\item $v$：车手的速度，单位 m/s。
\item $v_w$：风速，单位 m/s。
\item $v_r$：车手与风的相对速度，单位 m/s。
\item $\phi$：$v$ 的方向与 $v_w$ 的方向之间的夹角，即风向的方位角，单位 rad。
\item $P_{eff}$：驱动力的有效功率，即扣除传动损失后真正用于推动自行车前进的功率，单位 W。
\item $\lambda$：车手功率输出的转换效率，即有效功率与瞬时功率之比。
\item $\mu$：有效摩擦系数，用于描述滚动摩擦与传动摩擦的合成效应。
\item $\mu_s$：静摩擦系数，用于判定过弯时轮胎是否打滑。
\item $F_N$：轮胎与地面之间的法向力，单位 N。
\item $E_p$：势能，单位 J。
\item $E_k$：动能，单位 J。
\item $m$：车手与自行车的总质量，单位 kg。
\item $\theta$：坡道的倾角，上坡为正、下坡为负，单位 rad。
\item $s$：车手自起点起已经覆盖的距离，单位 m。
\item $R$：曲率半径。
\item $v_{safe}$：安全速度，即保证过弯不打滑的速度上限，单位 m/s。
\item $f_{\tau}$：切向摩擦力，单位 N。
\item $f_n$：法向摩擦力，即维持圆周运动所需的向心力，单位 N。
\end{itemize}

\section{分析与建模}

\subsection{模型一：功率输出与恢复}

\subsubsection{功率曲线}

从能量消耗的类型来看，人体有两种工作状态：有氧运动与无氧运动。在有氧运动过程中，人体消耗由有氧系统、糖原—乳酸系统以及磷酸原系统所产生的 ATP 中的能量。尽管有氧呼吸在运动过程中能够持续产生能量，但其功率是有限的。当所需功率超过有氧运动的功率上限时（该上限定义为临界功率 $P_c$），人体将转而消耗无氧能源中的能量。然而，在持续高于临界功率的状态下工作时，人体能够利用的无氧能量总量是有限的。用于衡量这种耐力能力的能量总量称为无氧工作能力（Anaerobic Work Capacity，AWC）。换言之，$P_c$ 刻画的是“功率的可持续水平”，而 $AWC$ 刻画的是“超出该水平之后还能坚持多久”，两者共同决定了一名车手的功率—时间特征。

只要功率输出低于临界功率，有氧运动就可以无限地持续下去。而无氧运动只能维持有限的时间，之后当车手以低于临界功率的功率骑行时便进入恢复阶段，以重新获得足够的无氧能量。因此，我们把车手的运动阶段划分为恢复阶段（有氧运动）与工作阶段（无氧运动），并在此基础上进行如下分析。这一划分使得整条赛道上的骑行被化为一串交替出现的充能与耗能区间，从而为后续的逐路段推进提供了统一的描述框架。

记 $P$ 为瞬时功率，$w$ 为无氧能量。在工作阶段，无氧能量的变化率 $\dfrac{dw}{dt}$ 等于临界功率与瞬时功率之差：
\[
\frac{dw}{dt}=P_c-P
\]
该式表明，当车手以高于临界功率的功率骑行时，无氧能量以两者的差值速率被消耗；当功率低于临界功率时，无氧能量则以相同的速率得到回补。根据 Morton, R.H. 与 Billat, L.V.[13] 的工作，在恢复阶段该关系依然成立，只是 $\dfrac{dw}{dt}$ 的取值为正。也就是说，耗能与充能可以被同一个微分方程统一描述，只是功率相对于临界功率的高低不同。

现在，我们试图通过给出 $P$--$T$ 函数来确定一条功率曲线，用以表征车手持续维持某一功率输出水平的能力，其中 $P$ 为给定的功率，$T$ 为其能够持续的时间。这样一条曲线是全部后续分析的出发点：它把抽象的生理能力转化为一条可以直接读数的边界，任何超过该边界的功率安排在实际中都无法被维持。

考虑一个以恒定功率 $P$ 进行的过程。无氧呼吸中的能量消耗速率为瞬时功率与临界功率之差：
\[
w(t)=w(0)-(P-P_c)t
\]
该式由前一式直接积分得到，其中 $w(0)$ 为过程开始时的无氧能量余额。假定车手能够充分利用剩余的无氧工作能力，则当 $t=T$ 时能量被恰好耗尽，即 $w(T)=0$，于是可得方程：
\[
P=P_c+\frac{w(0)}{T}
\]
当 $T$ 趋于 0 时，$P$ 趋于无穷大。在这一假设下，车手将具有无限的瞬时爆发功率，这显然不合理。因此，我们认定车手的功率输出存在一个上限。这一修正的必要性在于：人体在极短时间内的功率输出受肌肉收缩速度、神经募集能力等硬性条件限制，不可能随着持续时间的缩短而无界增长。

设瞬时功率的上限为 $P_{lim}$。根据文献，做功能力与无氧能量呈负相关，即随着无氧运动的进行，$P$ 会逐渐下降。因此引入系数 $\alpha$ 来描述最大功率与无氧能量之间的关系：
\[
P_{lim}=\alpha w(0)+P_c
\]
在这一假设下，当 $T=0$ 时 $P_m=\alpha \cdot AWC+P_c$ 是一个有限值，这与实际情况相符。$\alpha$ 的物理含义是每单位无氧能量所能支撑的额外功率，因此其量纲为时间的倒数；$\alpha$ 越大，车手在能量充裕时的瞬时爆发越强，但能量耗尽后的衰减也越剧烈。

当车手刚进入工作阶段时 $P_{lim}>P$，功率尚可维持，直到 $t=T$ 时两者相等而达到临界状态。当 $t>T$ 时 $P_{lim}<P$，功率便无法维持。基于上述分析，我们可以求出 $T$：令 $P=P_{lim}$、$t=T$、$w(0)=AWC$，得
\[
P=\alpha\left[AWC-(P-P_c)T\right]+P_c
\]
该式把“功率上限随时间下降到与所要求功率相等”这一临界条件写成了代数方程，解出 $P$ 与 $T$ 之间的关系即得
\[
P=\frac{\alpha AWC+\alpha P_c T+P_c}{1+\alpha T}
\]
这一函数具有两条重要的极限性质：当 $T\to 0$ 时 $P\to \alpha AWC+P_c$，即功率上限是一个有限值；当 $T\to\infty$ 时 $P\to P_c$，即长时间骑行所能维持的功率趋近于临界功率。这两条性质恰好刻画出功率曲线的两端行为，中间部分则表现为随持续时间增加而单调下降的过渡段。

对于给定的车手，只要测定参数 $AWC$、$\alpha$、$P_c$，即可得到其功率曲线。基于 J. Pinot 与 F. Grappe（2011）论文中运动员的生理数据，对模型中的方程进行拟合可得到最佳拟合结果，即 $\alpha=0.0399\ \mathrm{s^{-1}}$，$P_c=354.5\ \mathrm{W}$，$AWC=24010.41\ \mathrm{J}$。拟合所得的曲线即为该车手的功率曲线。

拟合得到的功率—持续时间函数为：
\[
P(T)=\frac{316.4T+32880}{T+25.02}
\]
其中 $P$ 的单位为瓦特，$T$ 的单位为秒。可以看到，该函数的分子与分母均为 $T$ 的一次函数，因此其极限值与前述分析一致：$T$ 很小时功率约为 1314 W，$T$ 很大时功率趋近于 316.4 W。

\subsubsection{三类典型车手的功率曲线}

接下来，我们考察不同类型车手的功率曲线。根据文献数据，计时赛专家的临界功率（380.2 W）与无氧工作能力（26143 J）都超出常规水平，这使得该车手具有更强的耐力以及应对多种复杂地形的能力；而冲刺手的 $\alpha$ 值（0.043）远超车手的平均水平，这意味着其 $P_{lim}$ 相当大，具有惊人的爆发力；相对而言，女性车手的临界功率、无氧工作能力与 $\alpha$ 均小于男性车手，分别为 321.3 W、22189 J 与 0.035。

把上述三组参数分别代入功率曲线表达式，即可得到三类车手的曲线。三者的差异集中体现在曲线的形状上：计时赛专家的曲线在中长时间段上维持在较高水平，下降较为平缓，说明其更擅长在整条赛道上稳定地输出功率；冲刺手的曲线在极短时间内可达很高的峰值，但随持续时间增加衰减迅速，说明其优势区间极短；女性车手的曲线整体位于前两者之下，其临界功率与无氧工作能力均较小，因此其可持续功率水平与总可支配能量都相对有限。这一比较说明，模型确实能够用少量可测参数区分不同类型的车手，从而满足“对不同运动员具有普适性”的要求。

\subsection{模型二：最优功率分配}

\subsubsection{功率与能量分析}

在模型一中，骑行过程中的输出功率存在一个上限，表示为 $P_{lim}=\alpha w+P_c$。然而在实际情形中，车手为了取得更好的成绩可以突破该上限。根据生理学以及题目的要求，突破上限会给车手带来额外的负担，也就是说，车手的最大功率会相应降低。为了描述这部分超额功率对车手体力的影响，我们定义 $W_{ex}$ 为车手在功率上限之上所做的功。车手突破上限所做的额外功会导致系数 $\alpha$ 衰减，进而影响其无氧运动能力：
\[
\alpha(W_{ex})=\alpha(0)-0.00001\alpha W_{ex}
\]
其中 $W_{ex}$ 的单位为 J。该式表明，超额做功越多，$\alpha$ 衰减得越快，车手在随后阶段能够调用的瞬时功率上限也就越低。这一项惩罚在优化中起到约束作用：它使得“不顾一切地提前透支”不再是无代价的选择，从而让模型给出的解更接近真实的骑行行为。

\subsubsection{天气条件分析}

自行车比赛中的天气条件指的是风的条件，包括风速与风向。为简化计算，在本模型中我们假定风速与风向保持恒定。这一处理的合理性在于：在一条计时赛赛道上，风场随时间的变化相对于整场比赛的持续时间而言较为缓慢，因此把风近似为常量可以在保留主要效应的同时避免引入难以标定的时变参数。

设空气阻力为 $F_{air}$。需要注意的是，$F_{air}$ 并不一定阻碍车手前进。例如，当风向与车手运动速度的方向相同时，风可能起到加速作用。为了描述该力的方向，我们采用矢量分析。定义 $v$ 为车手的速度，$v_w$ 为风速。根据空气动力学，风阻可表示为
\[
\vec{F}=-\frac{1}{2}\rho|\vec{v}_r|^2 C_D A\cdot\frac{\vec{v}_r}{|\vec{v}_r|}
\]
其中 $\rho$ 为空气密度，$v_r$ 表示车手与风的相对速度，$C_D$ 与 $A$ 分别表示阻力系数以及车手与自行车的迎风面积。为简化表达，定义 $S_{eff}$ 为 $C_D A$。该式中的负号表示阻力方向与相对速度方向相反，而末尾的单位矢量则负责把大小与方向分离，使表达式在顺风与逆风两种情形下都自动取得正确的符号。

为了把矢量表达式量化，我们建立平面直角坐标系，取东为 $x$ 正方向、北为 $y$ 正方向。设风速 $\vec{v}_w$ 的方位角为 $\phi$，运动方向的方位角为 $\gamma$，则可列出如下表达式：
\[
\begin{cases}
\vec{v}_w=|\vec{v}_w|(\cos\phi,\ \sin\phi)\\[0.3em]
\vec{v}=|\vec{v}|(\cos\gamma,\ \sin\gamma)\\[0.3em]
\vec{v}_r=\vec{v}-\vec{v}_w=(v\cos\gamma-v_w\cos\phi,\ v\sin\gamma-v_w\sin\phi)
\end{cases}
\]
风阻仅沿切向对车手产生影响。设风对车手的切向力为 $F_{air}$，则风对车手的有效作用力为
\[
F_{air}=-\frac{1}{2}\rho|\vec{v}_r|^2 S_{eff}\cdot\frac{\vec{v}_r}{|\vec{v}_r|}\cdot\frac{\vec{v}}{|\vec{v}|}
\]
当 $F_{air}>0$ 时，风对车手起阻碍作用；当 $F_{air}<0$ 时，风对车手起推动作用。也就是说，同一阵风在赛道的不同路段上可能表现为阻力，也可能表现为推力，这正是赛道走向与风向之间的相对关系所决定的。

\subsubsection{动力学分析}

把车手与自行车视为一个整体。作用在该整体上的力包括车手的有效驱动力 $F_{drive}$、空气阻力 $F_{air}$ 以及内外部摩擦力，它们共同导致动能与势能的变化。受力分析示意中给出了上述各力的方向与作用关系。需要注意的是，在示意图中 $F_{drive}$ 是一个假想的等效驱动力，$F_{eff}$ 是有效摩擦力，二者分别代表功率输入与机械损耗在整车层面上的等价作用。

定义有效输出功率为 $P_{eff}=\lambda P$，其中 $\lambda$ 是描述车手功率有多大比例能够用于驱动自行车的系数。根据能量变化约束，可列出如下方程：
\[
P_{eff}-\frac{1}{2}\rho v_r^2 S_{eff}\cdot v-\mu F_N\cdot v-\frac{\Delta E_p}{\Delta t}-\frac{\Delta E_k}{\Delta t}=0
\]
其中 $\mu$ 为有效摩擦系数，$F_N$ 表示作用在轮胎上的法向力，$E_p$ 表示势能，$E_k$ 表示动能。该式本质上是一个能量守恒关系：车手输入的有效功率，一部分被空气阻力耗散，一部分被摩擦耗散，其余部分则转化为势能与动能的变化率。

对于给定的赛道，我们把距离划分为大量路段（在我们的模型中共有 1164 个路段）。假设每一个路段都是倾角相同为 $\theta_i$ 的坡道，上坡时 $\theta>0$，下坡时 $\theta<0$；同时我们采用等距分段，即每一段在水平方向上具有相同的距离。记第 $i$ 段的速度为 $v_i$，则可得如下关系：
\[
\begin{cases}
\dfrac{\Delta E_p}{\Delta t}=mg\sin\theta\\[0.6em]
\dfrac{\Delta E_k}{\Delta t}=mv\dfrac{\Delta v}{\Delta t}\\[0.6em]
F_N=mg\cos\theta
\end{cases}
\]
其中 $m$ 为车手与自行车的总重量，$g$ 为重力加速度。第一式给出重力做功引起的势能变化率，第二式给出速度变化引起的动能变化率，第三式给出轮胎所受的法向力。联立式（12）与式（13）并代入有效功率，可得
\[
\frac{\Delta v}{\Delta t}=\frac{\lambda P_i}{mv_i}-g\sin\theta_i-\mu g\cos\theta_i-\frac{\rho v_r^2 S_{eff}}{2m}
\]
该式即车手运动的加速度方程，其中右端四项依次对应驱动力、重力分量、摩擦阻力与空气阻力对加速度的贡献。

基于 JAMES C. MARTIN（2006）论文中的数据，我们取 $\lambda=97.7\%$，$\mu=0.0025$。不同姿态下骑行的车手具有不同的有效迎风面积：以站立姿态（工作阶段）骑行的车手 $S_{eff}=0.304\ \mathrm{m^2}$；以坐姿（恢复阶段）骑行的车手 $S_{eff}=0.245\ \mathrm{m^2}$；空气密度在 293.15 K、1 atm 下计算为 $1.205\ \mathrm{kg/m^3}$。把这些数据代入式（14），可以得到 $v$ 与 $t$ 的如下关系式：
\[
\frac{dv}{dt}=\frac{a}{v}+bv^2+c
\]
其中 $a$、$b$、$c$ 的取值已知。这三项分别对应功率项、空气阻力项与重力—摩擦项，因此该方程在结构上把能量输入与各类损耗清晰地分离开来。对式（15）作代数变形，利用 $\dfrac{ds}{dt}=v$，得到如下微分方程：
\[
\frac{dv}{ds}\frac{ds}{dt}=\frac{dv}{ds}v=\frac{a}{v}+bv^2+c
\]
其中 $s$ 表示车手自起点起已经覆盖的距离。这一变形把以时间为自变量的方程改写为以距离为自变量的方程，从而使我们能够在给定路段长度的情况下直接求解出口速度，而不必逐时刻推进。这是一个关于 $v$ 与 $s$ 的微分方程，其边界条件为 $s=S_i$ 时 $v=v_i$。通过求解该微分方程，可以得到 $v$ 关于 $s$ 的函数；对于划分出的第 $i$ 个路段，第 $i$ 个节点的距离 $S_{i+1}$ 已知，代入 $s=S_{i+1}$ 即可得到 $v_{i+1}$ 的数值解。

至此我们可以求解 $v_{i+1}$。对该微分方程进行数值积分，可得
\[
\int_0^{t_i}dt=\int_{v_i}^{v_{i+1}}\frac{dv}{\dfrac{a}{v}+bv^2+c}
\]
从而可以解出第 $i$ 个路段所用的时间 $t_i$。在此基础上，我们可以得到比赛的总时间 $T_{total}=\sum_{i=1}^{1164}t_i$。由于每一个路段的出口速度同时又是下一路段的入口速度，这种逐段推进的方式自然地把整条赛道串联为一个整体，因而任何一处功率的改变都会通过速度的传递影响其后所有路段的时间。

上述讨论对赛道上的任意地形都成立。对于赛道中的弯道，车手在转弯过程中的运动可以近似为圆周运动。轮胎所受的摩擦力沿切向与法向两个方向作用。切向摩擦力 $f_{\tau}$ 即上文讨论的阻力，而法向摩擦力 $f_n$ 是维持圆周运动的向心力（相对于摩擦力而言，风阻的法向分量可以忽略不计）。

列出圆周运动的动力学方程：
\[
f_n=\frac{mv^2}{R}
\]
其中 $v$ 为自行车的速度，$R$ 为弯道处的曲率半径。我们已将赛道划分为 1164 个路段。对于每一个倾角为 $\theta_i$、横截面半径为 $r_i$ 的路段，可以将其近似为一条等距螺旋线，该螺旋线的等效曲率半径为
\[
R_i=\frac{r}{\sqrt{1+\tan^2\theta_i}}
\]
地面与轮胎之间存在摩擦力的最大值，即静摩擦力 $\mu_s mg$，其中 $\mu_s$ 为轮胎与地面之间的静摩擦系数。为保证安全，我们规定自行车在转弯过程中不得打滑，即切向摩擦与法向摩擦的合力不得超过静摩擦力的上限：
\[
f_{total}=\sqrt{f_n^2+f_{\tau}^2}\le\mu_s mg
\]
联立式（17）与式（19），可得第 $i$ 段速度的上限 $v_{i,safe}$：
\[
v_i\le v_{safe}=\sqrt[4]{\left(\mu_s gR_i\right)^2-\left(\frac{f_{i,\tau}R_i}{m}\right)^2}
\]
这一上限在优化中充当硬约束：无论功率分配方案在时间上多么有利，只要某一路段的速度超过 $v_{safe}$，该方案便被判定为不可行。

\subsubsection{功率分配的优化}

对车手策略的优化，本质上是在所有路段上寻找使总时间最小的最佳功率分配的过程。为了得到全局最优解而不是局部最优解，我们选择模拟退火算法（Simulated Annealing Algorithm，SA）。该算法允许在搜索过程中以一定概率接受使目标函数变差的解，从而能够跳出局部极小，这一点对本题尤为重要，因为赛道上的功率分配具有上千个自由度，目标函数曲面存在大量局部极小。结合题目要求与常识，该优化问题中存在动力学约束、速度约束、能量约束与功率约束。优化模型表述如下：
\[
\begin{aligned}
\min\ & J=T_{total}=\sum_{i=1}^{1164}t_i\\
\text{s.t.}\ &
\begin{cases}
\text{动力学模型：}\ \dot{s}=f(s,P)\\
\text{剩余能量限制：}\ 0\le w_i\le AWC\\
\text{速度限制：}\ 0\le v_i\le v_{i,safe}\\
\text{总能量限制：}\ \displaystyle\int_0^{T_{total}}dt\le E_{lim}
\end{cases}
\end{aligned}
\]
在建立目标函数与约束条件之后，我们的模拟退火算法的思路如下：

\begin{itemize}
\item \textbf{第 0 步：}设定决策变量为 $P_i\ (1\le i\le 1164)$，初始解为 $P_i=P_c\ (1\le i\le 1164)$，描述算法进程的参数为温度 $K$，退火函数为 $K_{j+1}=aK_j$。确定初始温度与终止温度的取值。以临界功率作为初始解，意味着搜索从“可以无限维持的恒定功率”这一保守方案出发，此后的一切改进都以不违背生理约束为前提。

\item \textbf{第 1 步：}依次对 100 个路段施加随机扰动，扰动量 $\Delta P$ 服从均值为 0 的正态分布。每次扰动之后，我们检查约束是否满足：若约束满足，则保留该扰动；若约束不满足，则不保留该扰动。每一次扰动的作用方式在示意图中给出。这种“先扰动、后校验”的做法保证了搜索始终在可行域内进行，避免了不可行解对后续迭代的污染。

\item \textbf{第 2 步：}在完成 100 次扰动之后，计算该情形下完成比赛所花费的总时间 $f_1$，并与旧解 $f_0$ 进行比较：当新解更小时，选取新解；当旧解更小时，以一定概率选取新解。选取新解的概率为
\[
P=
\begin{cases}
1, & f_1<f_0\\[0.6em]
e^{-\frac{f_1-f_0}{T}}, & f_1>f_0
\end{cases}
\]
该准则表明，目标函数变差时被接受的概率随温度 $T$ 的降低而迅速减小：在搜索初期温度较高，算法愿意接受较差的解以扩大搜索范围；在搜索后期温度较低，算法几乎只接受改进解，从而在最优解附近细化。

\item 完成一次迭代后，用退火函数模拟系统“温度”的下降，并判断是否达到终止温度 $K_f$。

\item \textbf{第 3 步：}重复第 1 步与第 2 步，直到“温度”达到终止温度。
\end{itemize}

上述迭代流程以一个完整的循环为单位不断重复，每次循环都在当前温度下对功率分配作局部试探，并按 Metropolis 准则决定是否接受，从而在搜索初期允许目标函数暂时变差以跳出局部极小，而在温度逐步降低后逐渐收敛到全局最优解附近。由于每次扰动只涉及 100 个路段，单次迭代的计算量被控制在可接受范围内，而多次迭代累积起来则足以覆盖整条赛道的全部路段。
\subsection{2021 年奥运会个人计时赛赛道}

东京奥运会赛道以其陡峭的坡度而著称，这要求车手灵活地调整输出功率，以利用下坡过程中的加速过程。也就是说，在这条赛道上，功率分配的意义不仅在于“省力”，更在于“借势”：车手在下坡时不应当把速度的增长视为理所当然，而应当把它视为一段可以免费获得的速度增量，从而在上坡之前主动压低功率以储备能量，并在下坡中把储备转化为动量。

首先，我们假定车手具有平均的身体能力。若他以恒定的临界功率 $P_c$ 输出功率，则其完成比赛所花费的总时间为 6146 s。这一恒功率方案可以看作所有策略的基准线：它的优点是无氧能量几乎不被动用，因而在任何时刻都不存在能量耗尽的风险，但它的缺点同样明显，即车手在全程都没有利用无氧工作能力所提供的额外速度空间。本模型所设计的模拟退火算法重新分配了车手在整条赛道上的功率，以有效利用无氧工作能力，所得结果如下。

从时间上看，车手的成绩提升了 29\%。由于模拟退火算法本身具有随机性，功率曲线中存在一定的噪声，因此我们对其作了平滑处理。平滑仅仅去除了相邻路段之间高频的功率抖动，并未改变功率的整体走势，因此优化所得的节律依然清晰可辨。

无氧工作能力的分布与速度的分布呈现出高度一致的对应关系：在功率被抬高的路段，无氧能量被持续消耗，速度随之上升；而在功率被压低的路段，无氧能量重新累积，速度则相应下降。由此可以看出，最优方案并不是让无氧能量始终保持在满值，也不是让它一路耗尽，而是让它在赛道上有节奏地“充—放”，从而在全程范围内被均匀地使用。

在模拟退火算法的作用下，模型给出了若干条更为合理的策略。其一是在地形较为平缓处让功率在高与低之间交替，以在保持速度的同时确保在上坡之前不过度消耗能量；其二是在上坡时提高功率、在下坡时降低功率以恢复体力；其三是在接近终点处全力冲刺，把剩余的体力全部消耗掉。这三条策略共同构成了一条完整的比赛逻辑：前半程控制消耗、中段借地形节律调节、末段清空储备。

优化过程的曲线表明，我们的模型在优化方面表现良好：目标函数（总时间）随着迭代次数的增加而迅速下降，随后进入缓慢的细化阶段，并在退火温度降低到终止温度附近时基本稳定。这一形状说明算法已经从初始的恒功率解出发，稳定地找到了明显更优的功率分配方案，且没有在迭代末期出现明显的反弹，即所得结果不是搜索过程中的偶然波动。

从物理机制上看，29\% 的改进幅度并不神秘。以临界功率恒速骑行时，整场比赛的无氧工作能力几乎完全未被使用，车手相当于放弃了一块可以自由动用的“速度预算”；而优化后的方案把这块预算按赛道地形的需要分期支出，在爬坡与加速阶段集中投放，在下坡与缓段回收。由于总时间是速度倒数的积分，在低速段上增加同样的功率所带来的时间缩短，远大于在高速段上增加同样的功率所带来的时间缩短，因此把额外的功率投放在速度最低的上坡段，其边际收益远高于把它平均分摊到全程，这正是优化能够显著缩短总时间的根本原因。反过来，若赛道上不存在明显的低速段，同样的预算所能带来的边际收益便大打折扣，改进幅度也随之下降。

\subsection{2021 年 UCI 世界锦标赛个人计时赛赛道}

我们将把同样的研究方法应用于 UCI 赛道。由于 UCI 赛道相对平坦，其对功率分配的要求不如第一条赛道那样高。策略基本上是在高功率与低功率之间交替。成绩的提升也不算显著，为 9\%。

这一结果本身具有解释价值：当赛道几乎没有长距离爬坡时，重力项对速度的影响被削弱，车手在全程所面对的主要阻力是空气阻力，而空气阻力只与速度有关、与所在位置无关，因此“在何处用力”的收益远小于第一条赛道。此时最优方案的主要作用变成了维持一个稳定的高速巡航，并通过小幅的功率交替来对抗风阻与微小起伏。

该赛道的功率分配曲线与优化过程曲线同样表明，算法在较平坦的赛道上仍然能够收敛，但由于可挖掘的空间有限，功率曲线的起伏幅度明显小于第一条赛道，优化的收益也随之下降。从优化过程曲线还可以看到，算法在这条赛道上收敛得更快，前若干次迭代即已取得大部分收益，此后目标函数的下降十分平缓，这与赛道上缺乏决定性难点的事实相一致：当各处的地形条件彼此接近时，功率分配的可优化维度也随之减少。

\subsection{自行设计的赛道}

在我们自行设计的赛道中，我们布置了四个急弯与一段陡坡，地形自然的凹凸起伏由一系列服从正态分布的点经过三次样条插值得到。之所以采用样条插值，是为了保证生成的地形曲线连续且光滑，从而使坡度不会出现物理上不合理的突变，同时又能保留足够的随机性以模拟真实赛道的复杂性。

我们的算法在这条赛道上的效果介于第一条赛道与第二条赛道之间，与恒功率骑行相比提升了 18\%。该赛道的功率分配情况表明，在急弯之前功率被适当压低以降低入弯速度、保证过弯安全，在陡坡段功率被显著抬高，而在长下坡段则出现明显的功率低谷以恢复体力。值得注意的是，急弯处的功率下调并非出于省力考虑，而是由安全速度约束所强制的：当曲率半径较小而速度较高时，所需的向心力会迅速增大并逼近静摩擦上限，模型因此自动把速度压回可行范围之内。

设置这条自行设计的赛道，目的在于检验模型在“两种特征同时出现”时的表现：两条真实赛道分别以陡坡密集与平坦单向为主要特点，而自行设计的赛道则把急弯与陡坡压缩在较短的赛程之内，使约束之间的冲突更为尖锐。在这种情形下，模型必须在“提高功率以爬坡”与“压低速度以过弯”之间反复权衡，其结果既不是前两条赛道结果的简单叠加，也不是任何一种单一节律的重复，这说明模型对复合地形具有独立的适应能力。

从这三条赛道可以看出，当坡度较大、地形不平时，我们的数学模型与优化过程具有很高的优越性。反之，当地形平坦且缺乏急弯时，模型的收益主要体现在对能量节奏的精细调配上，其改进幅度相对有限。这一对比也提示我们，模型的价值并不在于对所有赛道都给出同样惊人的提升，而在于它能够针对具体赛道定量地指出改进空间究竟有多大、改进应当发生在哪些路段。

\section{敏感性分析}

\subsection{天气条件}

我们首先对天气条件进行敏感性分析。在模型二中，我们分析了空气阻力对车手作用的力学机理，并把有效作用力 $\vec{F}_{air}$ 表示为 $\vec{v}_r$ 与 $\vec{v}$ 的函数。因此，天气条件对结果的影响完全是通过相对速度 $\vec{v}_r$ 这一中间量传递的：风速与风向改变的是 $\vec{v}_r$ 的大小与方向，而 $\vec{v}_r$ 又决定了空气阻力的大小与正负。

为便于建模，我们设定风速与风向服从具有给定数学期望的正态分布。这样一来，每一组期望值都对应一个可重复的“天气情形”，而正态分布本身则刻画了真实风场中不可避免的随机涨落，使分析结果不至于过分依赖于某一个精确的风速取值。

接下来，我们以 2021 年奥运会个人计时赛赛道（赛道一）与 2021 年 UCI 世界锦标赛个人计时赛赛道（赛道二）为例，说明不同风向与风力大小对结果的影响。我们之所以选择这两条赛道，是因为它们的整体走向恰好构成对照：赛道一总体呈环形，各方向的路段长度大致均衡；赛道二总体由北向南蜿蜒，具有明确的主方向。

我们继续沿用模型二中的平面坐标系，把风向的方位角记为 $\phi$。

我们首先研究风向固定的情形，其中 $\phi\sim N(\pi,0.15)$。令风速期望 $E_{v_w}$ 由小到大变化，考察该风速下的优化结果。

从结果可以看出，对于赛道一，随着风速期望的增大，比赛成绩的优化结果在一个很小的范围内上下波动，与风速并未表现出显著的相关性。这说明当我们的模型应用于赛道一这类环形赛道时，其对风速的敏感性相当弱。其原因在于，环形赛道上顺风路段与逆风路段交替出现，风速增大所带来的收益与损失在很大程度上相互抵消，因而总时间对风速几乎不敏感。

对于赛道二，随着风速期望的增大，比赛成绩的优化结果出现明显的改善，这是因为赛道二总体上是由北向南蜿蜒的。因此，像赛道二这样的赛道往往对风速具有很强的敏感性。换句话说，当赛道的走向与风向大体一致时，风不再是一个可以在全程被平均掉的扰动，而成为一个系统性的助力或阻力。

随后我们固定风速期望 $E_{v_w}$，令风向期望 $E_\phi$ 在 $(0,2\pi)$ 内增大。可以看到，赛道一的优化结果只出现小幅的随机波动，这表明在赛道一这类环形赛道中，与其他类型的赛道相比，我们的模型对风向的敏感性较低。对于赛道二，其优化结果的波动与赛道的整体走向之间存在显著的相关性（顺风时用时更少，逆风时用时更多），且波动幅度更大，说明赛道二对风向具有很强的敏感性。

综合风速与风向两个方向的分析可知，赛道越接近环形，模型对天气条件的敏感程度越低；而赛道走向越单一、越具有明确的主方向，模型对天气条件就越敏感。这一结论具有直接的实战含义：在环形赛道上，车队不必为当天的风况投入过多的准备；而在单向型赛道上，风况信息应当被纳入赛前的功率规划之中。

需要补充说明的是，这里的敏感性并不是指模型对输入数据的数值精度敏感，而是指最优解的结构是否会被天气参数改变。在环形赛道上，风向改变时最优解的节律基本保持原有形状，只是各路段的具体功率数值略有平移；而在单向型赛道上，风向改变会使某一段原本应当“省力”的路段变成应当“发力”的路段，最优解的结构因此发生实质变化。这一区别解释了为什么两类赛道的敏感程度会出现量级上的差异，也提示我们在把模型推广到其他赛道时，应当先判断赛道属于哪一类，再决定天气参数需要以多高的精度被标定。

\subsection{对目标功率分配的偏离}

以赛道三为对象，本部分将对车手实际功率分配与目标功率分配之间的偏离程度进行敏感性分析。选择赛道三的原因是它同时包含急弯与陡坡，能够较全面地暴露执行偏差的影响。我们对车手的目标功率分配施加如下操作：
\[
\begin{cases}
P_i'=P_{target,i}+\epsilon_i\\
\epsilon_i\sim N(0,\epsilon)
\end{cases}
\]
其中 $\epsilon$ 是用于量化实际偏离程度的指标。该指标的含义是：车手在每个路段上的实际功率围绕目标功率作无偏的随机波动，波动幅度由 $\epsilon$ 控制。然而需要指出，经过上述操作得到的功率分配 $P_i'$ 可能并不满足我们前面提到的约束条件，例如它可能使无氧能量在某个时刻超出上限，或使某个路段的速度超过安全值。因此，还需要执行如下操作：

\begin{enumerate}
\item 当 $i=1$ 时，由初始条件与约束条件得到 $\Omega_1$，即 $P_1$ 的取值范围：
\[
P_{real,1}=
\begin{cases}
P_1', & P_1'\in\Omega_1\\
\max\{\Omega_1\}, & P_1'>\max\{\Omega_1\}\\
\min\{\Omega_1\}, & P_1'<\min\{\Omega_1\}
\end{cases}
\]
也就是说，凡超出可行范围的取值都被截断到边界上，从而保证起点处的状态在物理上是可实现的。进而得到 $t_{real,1}$ 与 $w_{real,1}$。

\item 当 $i=k\ (2\le k\le 1000)$ 时，根据 $w_{real,k-1}$ 与约束条件得到 $\Omega_k$：
\[
P_{real,k}=
\begin{cases}
P_k', & P_k'\in\Omega_k\\
\max\{\Omega_k\}, & P_k'>\max\{\Omega_k\}\\
\min\{\Omega_k\}, & P_k'<\min\{\Omega_k\}
\end{cases}
\]
进而得到 $t_{real,k}$ 与 $w_{real,k}$。可以看到，每一个路段的可取值范围都依赖于上一路段结束时的无氧能量余额，因此截断操作必须逐段顺序进行，而不能独立地对所有路段一次性完成。
\end{enumerate}

依次对每一个路段作上述调整，以保证功率分配满足约束条件，然后计算 $T_{real}=\sum_{i=1}^{1000}t_{real,i}$。逐步增大 $\epsilon$ 的取值，并分别计算相应的比值 $T_{real}/T_{target}$。

从结果可以看出，模型对车手功率分配的偏离是敏感的；并且随着偏离程度的增大，完成比赛所花费的总体时间趋于增加。这一正相关关系符合直觉：偏离越大，车手在部分路段所做的功就越偏离最优节律，而在另一些路段又被约束强制截断，因而总体上既浪费了能量又损失了速度。更值得注意的是，总时间的增长并非线性，而是在偏离达到某一水平后明显加速上升，这意味着执行精度的要求存在一个“临界带”。

本部分接下来首先分析关键路段。我们将对每一个单独的路段进行敏感性分析：对该路段上的功率施加 10\% 的扰动，敏感性较大的路段即被理解为关键路段。对第 $i$ 个路段的分析步骤如下。

\begin{enumerate}
\item 确定 $\Omega_i$，即 $P_i$ 的取值范围；
\item 令 $P_i(1)=1.1\cdot P_{target,i}$，$P_{real,i}(1)=\max\{P_i(1),\max\{\Omega_i\}\}$，即把该路段功率上调 10\% 并按约束截断；
\item 依次调整 $P_{real,i+1}(1),\ldots$，得到 $T_{real}(1)$；
\item 令 $P_i(2)=0.9\cdot P_{target,i}$，$P_{real,i}(2)=\min\{P_i(2),\min\{\Omega_i\}\}$，即把该路段功率下调 10\% 并按约束截断；
\item 依次调整 $P_{real,i+1}(2),\ldots$，得到 $T_{real}(2)$；
\item 于是得到灵敏度指标 $B_i=(T_{real}(1)+T_{real}(2))/T_{target}$。
\end{enumerate}

上述指标同时包含了上调与下调两种扰动的影响，因此它衡量的是该路段对功率误差的双向敏感程度，而不是仅对某一个方向的敏感程度。将所得的各路段灵敏度与赛道的地形图进行对比即可作出判断。

结论很容易分析：上坡路段（尤其是陡坡）对偏离程度的敏感性非常强；急弯的敏感性次之。其机理在于，上坡时重力分量与驱动力方向相反，功率的每一单位变化都会直接反映为速度的变化，而且速度一旦下降便需要更长的时间才能恢复；急弯处则由于安全速度约束的存在，功率的微小变化可能使速度触及上限或跌破下限，从而使该路段的通过时间发生跳变。据此，我们可以把 11$\sim$13 km 处的陡坡视为关键路段，把三个急弯前后各 100 m 的路段视为次关键路段。

我们规定，关键路段上的功率偏差所引起的成绩偏差不得超过 0.3\%，由此分别反向求解各关键部位可能的分段计时区间。这一规定的意义在于把“功率执行精度”这一难以直接观测的要求，转换为“分段用时”这一在比赛中可以实时读取的指标，从而使指导具有可操作性。

\begin{itemize}
\item 11$\sim$13 km 路段的通过时间区间为 $(337\sim401)$ s；
\item 9.9$\sim$10.1 km 路段的通过时间区间为 $(36.2\sim45.1)$ s；
\item 19.9$\sim$20.1 km 路段的通过时间区间为 $(37.4\sim43.2)$ s；
\item 29.9$\sim$30.1 km 路段的通过时间区间为 $(35.1\sim41.6)$ s。
\end{itemize}

上述四个区间可以直接作为赛中的检查点使用：只要车手在通过这些路段时用时落在相应区间之内，就可以认为其功率执行精度处于可接受的范围内；若用时明显偏离区间，则说明在该路段上出现了系统性的功率偏差，需要在后续路段中作出补偿。由于这些区间是由“成绩偏差不超过 0.3\%”这一要求反解得到的，它们同时也给出了执行精度的量化标准，而不必依赖车手对自身功率的主观感知。

\section{模型的推广}

团体计时赛的动力学模型如下：一组自行车手在骑行过程中排成一条直线，由于前排车手的破风作用，后排车手所受的空气阻力系数相对于单人骑行时有明显降低。因此，与个人计时赛相比，在团体计时赛中，车队应当选择恰当的策略轮流破风，这将成为提升车队成绩的关键。需要强调的是，团体计时赛的计分规则使得这一问题的性质发生了变化：成绩并不取决于第一名车手，而取决于第四名车手，因此车队的目标不是让某一个人最快，而是让“第四快”尽可能快。

\subsection{简化情形}

为简便起见，我们首先从一个不考虑队员之间不同分工的模型出发。在这一模型中，六名车手被视为六个可以互换位置的相同个体，唯一的决策是在每一个路段上采用哪一种排列。

根据文献，在比赛过程中，两辆自行车之间的距离越小，后排所受的风阻降低效果越明显。然而，为保证车手的安全，两名车手之间应当设定一个安全距离 $d_{safe}$；若距离小于 $d_{safe}$，就会存在追尾的风险。理想情况下，车手始终能够把两人之间的距离恰好保持在安全距离上，从而使平均风阻最小。由此我们可以假定车手的速度与加速度始终保持相同，下面的讨论均以此为基础。这一假定把原本需要分别追踪六个速度的复杂问题，简化为只需追踪一个速度的问题。

在一个六人车队中，容易知道共有 $6!=720$ 种排列。记这些排列为 $A_k\ (1\le k\le 720)$，并令
\[
A_k=
\begin{cases}
1, & \text{若选用排列 }A_k\\
0, & \text{若未选用排列 }A_k
\end{cases}
\]
在我们的模型中，赛道被划分为 1000 个路段，在每一个路段内排列保持不变。也就是说，队形的调整只在路段之间进行，这与实际比赛中车手在直道或缓坡上完成换位的做法相符。

我们沿用与个人计时赛相同的动力学模型。但在团体计时赛中，处于车队不同位置的车手所受的空气阻力 $F_{air}$ 并不相同，因为 $F_{air}$ 是 $v_r^2$ 的线性函数。设 $F_{air}=Cv_r^2$。令 $A_{ki}$ 表示第 $i$ 个路段上选用的排列 $k$，$c_{jk}$ 表示第 $k$ 种排列中第 $j$ 位车手的系数，则第 $j$ 位车手的系数为
\[
C_{j,i}=\begin{bmatrix}A_{1i},\cdots,A_{720i}\end{bmatrix}
\begin{bmatrix}c_{j,1},\cdots,c_{j,720}\end{bmatrix}^{T}
\]
该式的作用是把“当前选用了哪一种排列”这一离散决策，转化为每一位车手实际承受的阻力系数：由于 $A_{ki}$ 中只有一个分量为 1、其余为 0，上述内积恰好从 720 个候选系数中选出与当前排列对应的那一个。利用模型二中的方法，即可得到第 $i$ 个路段的时间 $t_i$。

\subsection{复杂情形}

考虑到车队成绩由第四位通过终点的车手成绩决定，教练需要为车手安排不同的职责。两名破风手需要在比赛的前半程承担更多的破风任务，以便其余四名车手为最后路段中的冲刺保存体力。因此，在何时让这两名车手脱离车队成为一个关键的选择：抛弃得过早，则剩余四人在长距离上失去遮挡，风阻骤增；抛弃得过晚，则两名破风手白白消耗体力，而四人组的冲刺储备不足。

设 $i_d$ 为车队抛弃破风手所在的路段。在第 $i_d$ 个路段之后，我们只需要考虑一个四人车队，即共有 $4!=24$ 种排列。在这种情况下，第 $j$ 位车手的系数可以表示为
\[
C_{j,i}=\begin{bmatrix}A_{1i},\cdots,A_{24i}\end{bmatrix}
\begin{bmatrix}c_{j,1},\cdots,c_{j,24}\end{bmatrix}^{T}
\]
与模型二相同，我们仍把它视为一个优化问题，并使用模拟退火算法求解。决策变量为 $i_d$、$P_i$ 的分配以及 $A_i$。目标函数为
\[
\min\ J=\sum_{i=1}^{i_d}t_i+\sum_{i=i_d+1}^{1000}t_i
\]
由于决策变量数量庞大（同时包含连续型的功率分配与离散型的排列选择，以及抛弃点这一整数变量），直接对全部变量同时搜索的计算代价过高。因此我们采用如下方法：

\begin{itemize}
\item \textbf{第 1 步：}$i_d$ 均匀地取 10 个离散值。对每一个 $i_d$，确定初始解、降温函数、初始温度与终止温度。这一步的作用是把一个混合整数优化问题分解为 10 个规模较小的子问题。
\item \textbf{第 2 步：}对第 1 步中的 10 个离散值，比较相应的最优目标函数值，即总时间。根据比较结果，取次峰值之间的区间，并把 $i_d$ 约束在该范围内。若区间内包含的点数仍然过多，则在该区间内重复第 1 步。其思想是先以粗网格定位最优点所在的大致区间，再在小区间内加密搜索，从而避免在全域上作穷举。
\item \textbf{第 3 步：}把 $i_d$ 加入决策变量，使用模拟退火算法得到最优方案。
\end{itemize}

这一分层求解的做法在实际中也有对应的赛场解释：教练组通常会先凭经验把“换人点”锁定在某几公里的范围内，再在该范围内细究具体位置与配速，而不是从一开始就同时决定全部细节。模型的求解流程与这一决策习惯一致，因此其结果更容易被教练接受和执行。

\section{讨论}

\subsection{模型的优缺点}

\begin{itemize}
\item 我们的动力学模型与实际情形吻合得较为紧密。模型同时考虑了重力分量、滚动摩擦、空气阻力与传动效率，并通过对赛道逐段离散的方式保留了地形与弯道的细节，因此其速度预测与时间预测具有明确的物理依据。
\item 模拟退火算法不能保证得到全局最优解。算法所得的方案在统计意义上是高质量的，但理论上仍存在更优的解未被访问到的可能，这一局限源自启发式算法本身的性质。
\item 每次对不同赛道进行优化时，模拟退火算法的参数都需要重新调整。初始温度、降温系数与终止温度对收敛速度与解的质量影响显著，而不同赛道的地形复杂度差异较大，因此参数难以一次设定后长期沿用。
\item 模型对所引用数据的准确性高度敏感。功率曲线中的生理参数、阻力系数与迎风面积均取自文献，这些数值的偏差会沿着模型逐段传播，并在总时间上被累积放大。
\end{itemize}

\subsection{未来工作}

\begin{itemize}
\item 收集更多数据，以提升模型的准确性与稳健性。特别是针对不同性别、不同专长的车手分别标定参数，可以使功率曲线的适用面进一步扩大。
\item 把模型应用于其他赛道，以检验模型的有效性。通过在不同类型赛道上的预测与实测对比，可以系统地评估模型的误差来源，并据此改进其中的简化假设。
\end{itemize}

\section{比赛指导}

尊敬的车队经理：

本指导将围绕一位参加日本东京 2021 年奥运会个人计时赛赛道的冲刺手的结果展开。

如您所知，您所负责的这位车手是一位典型的冲刺手。他具有非凡的短时爆发力，能够在短距离冲刺比赛中脱颖而出并轻松取胜。与计时赛专家相比，他的身体特征可以在功率曲线中得到很好的体现：在极短的持续时间内，他的功率上限明显更高；但随着持续时间的延长，其功率上限的衰减也更为迅速。这一特征意味着，他在这条赛道上的优势并不在于长时间的高速巡航，而在于关键时刻的短促加速。

然而，在即将到来的奥运会个人计时赛赛道中，复杂崎岖的地形以及大量的弯道，对运动员在长距离骑行中体力的合理分配与有效利用提出了更为严峻的挑战。

东京奥运会赛道以显著的坡度为特征，这要求运动员灵活地调整功率输出，并利用下坡的加速过程。如果车手没有妥善规划自己的功率输出，而是随意地分配体力，那么比赛结果对您和您的车队来说将会是令人失望的。对于这位以爆发力见长的车手而言，风险尤为突出：一旦他在前半程凭直觉反复冲刺，无氧能量便会在长坡到来之前被提前耗尽，此后即便地形有利，他也再无力加速。

幸运的是，我们的模型为您很好地解决了这一问题。

我们的模型把赛道划分为 1164 个等距路段。我们已经为您和您的车手提供了每一个路段的可靠推荐功率输出。总体而言，我们的模型建议您的车手在上坡时提高功率输出以提升速度，在下坡时降低功率以恢复体力，在平路上则让功率在高与低之间交替。如果严格按照我们的模型建议执行，您的车手的成绩最多可以提升 29\%，这一效果是惊人的。

然而，运动员在骑行过程中对自身功率输出的感知判断往往并不准确，这使得他们难以严格实施我们推荐的方案。因此，有必要在自行车上安装一套能够远程监测车手功率输出的传感器系统。您和您的团队可以利用耳机等工具，实时告知您的车手应当提高还是降低功率。这种“外部仪表 + 实时口令”的组合，可以把车手从不可靠的主观感觉中解放出来，使其专注于执行。

当然，尽管有传感器的辅助，一位在漫长的自行车骑行中经历剧烈无氧运动的车手，仍可能难以精确地控制自己的功率。在出现不可避免的偏差时，我们的模型已经识别出东京奥运会赛道上的关键路段。在这些关键路段上，您的车手需要极为严格地执行我们推荐的功率输出，否则可能对他的比赛成绩产生显著影响。这些关键路段包括：所有的陡坡与急弯。您的车手在通过这些路段时应当格外小心。在其余路段上，即使出现一定程度的执行误差，对总成绩的影响也相对有限，因此车手可以在那里适度依靠感觉骑行，以减轻注意力负担。

值得一提的是，您和您的车手不必过多担心当天的风速与风向。我们的空气动力学模型已经把风速与风向考虑在内，您的车手只需尽可能严格地遵循我们的建议即可。

\section*{参考文献}

\begin{itemize}
\item[{[1]}] https://olympics.com/olympic-games/tokyo-2020 \& TopographicMap
\item[{[2]}] https://www.flanders2021.com \& TopographicMap
\item[{[3]}] Vandewalle, H., 2018. Modelling of Running Performances: Comparisons of Power-Law, Hyperbolic, Logarithmic, and Exponential Models in Elite Endurance Runners. BioMed Research International, 2018, pp.1-23.
\item[{[4]}] Debraux, P., Grappe, F., Manolova, A.V. and Bertucci, W. (2011). Aerodynamic drag in cycling: methods of assessment. Sports Biomechanics, 10(3), pp.197218.
\item[{[5]}] MARTIN, J.C., GARDNER, A.S., BARRAS, M. and MARTIN, D.T. (2006). Modeling Sprint Cycling Using Field-Derived Parameters and Forward Integration. Medicine \& Science in Sports \& Exercise, 38(3), pp.592597.
\item[{[6]}] Blocken, B., van Druenen, T., Toparlar, Y., Malizia, F., Mannion, P., Andrianne, T., Marchal, T., Maas, G.-J. and Diepens, J. (2018). Aerodynamic drag in cycling pelotons: New insights by CFD simulation and wind tunnel testing. Journal of Wind Engineering and Industrial Aerodynamics, [online] 179, pp.319337.
\item[{[7]}] Pinot, J. and Grappe, F. (2011). The Record Power Profile to Assess Performance in Elite Cyclists. International Journal of Sports Medicine, 32(11), pp.839844.
\item[{[8]}] Peronnet, F. and Thibault, G. (1989). Mathematical analysis of running performance and world running records. Journal of Applied Physiology, 67(1), pp.453465.
\item[{[9]}] Pinot, J. and Grappe, F. (2011a). The Peronnet--Thibault mathematical model applied to the record power profile in cycling. Computer Methods in Biomechanics and Biomedical Engineering, 14(sup1), pp.219220.
\item[{[10]}] Vanhatalo, A., Jones, A.M. and Burnley, M. (2011). Application of Critical Power in Sport. International Journal of Sports Physiology and Performance, 6(1), pp.128136.
\item[{[11]}] Pinot, J. and Grappe, F. (2010). The Power Profile for determining the physical capacities of a cyclist. Computer Methods in Biomechanics and Biomedical Engineering, 13(sup1), pp.103104.
\item[{[12]}] Millet, G.P., Tronche, C. and Grappe, F. (2014). Accuracy of Indirect Estimation of Power Output From Uphill Performance in Cycling. International Journal of Sports Physiology and Performance, 9(5), pp.777782.
\item[{[13]}] Morton, R.H. and Billat, L.V. (2004). The critical power model for intermittent exercise. European Journal of Applied Physiology, 91(2-3), pp.303307.
\item[{[14]}] Blocken, B., Toparlar, Y., van Druenen, T. and Andrianne, T. (2018a). Aerodynamic drag in cycling team time trials. Journal of Wind Engineering and Industrial Aerodynamics, [online] 182, pp.128145.
\end{itemize}

\end{document}
